\section{Conclusions}
\label{sec:conclusions}

This study examines the energetic response of a lid-driven cavity to superelliptic bottom-wall shaping. For $n=2\to4$ at $\alpha=0$, fluid area increases by $7.93\%$, but integrated dissipation rises by only $0.0218\%$, $0.3249\%$, and $0.5909\%$ at $Re=100$, $400$, and $1000$ on the $N=384$ production grid. At $N=768$ the corresponding increases are $0.0223\%$, $0.3244\%$, and $0.5920\%$. An area-only reference overpredicts these contrasts by factors of $364$, $24$, and $13$. Dissipation per unit fluid area decreases by about $7\%$, showing that area change alone cannot explain the response.

An exact spatial partition separates dissipation in the newly opened shoulder region from redistribution over the region common to both geometries. At $Re=400$, the added region occupies $7.4\%$ of the fluid area but contributes $0.2\%$ of the reference dissipation, while the common-region redistribution is a near-cancellation of opposing lid-band and core-band contributions. The common-region contribution is negative at $Re=100$ and positive at $Re=400$ and $1000$; the net contrast remains positive throughout. The four bands were recomputed with the same exact-partition code at $N=128$, $256$, $384$, $512$ where available and $N=768$ at all three Reynolds numbers (Table~SM.XV, Fig.~SM.5). At $N=768$ the partition gives $\mathrm{d}A=-4.73$ and $\mathrm{d}B=+6.96$ at $Re=100$, $+12.62$ and $+19.82$ at $Re=400$, and $+9.99$ and $+49.21$ at $Re=1000$, in units of $10^{-4}D^*(2,0)$; the corresponding contrasts are $0.02232\%$, $0.32442\%$ and $0.59201\%$, within $2.5\%$, $0.14\%$ and $0.19\%$ of the $N=384$ values. The cancellation index rises from $4.0$ at $Re=100$ to $14.7$ at $Re=1000$, while the net contrast remains stable. The cell-by-cell identity $\Delta D=\mathrm{d}A+\mathrm{d}B$ is algebraically exact and closes to machine precision at every $N$. The numerical accuracy of the individual terms is weaker, because $\mathrm{d}A$ is the residual of larger opposing band contributions and its relative uncertainty grows with the cancellation index; the discrete partition as a whole is stable to a few percent, but the separate continuum limits of $\mathrm{d}A$ and $\mathrm{d}B$ are not established. The grid sensitivity of $\mathrm{d}A$ reaches $1.9\%$, $9.0\%$ and $30\%$ of its value at $Re=100$, $400$ and $1000$. At $Re=100$, the energy-balance defect varies by $5.0$ points between $N=128$ and $N=768$ while $I_D(2,4)$ remains within $7.2\times10^{-4}$~pp (Table~SM.XIX), showing that the defect cancels in the paired contrast.

The response is concentrated in the $n=2\to4$ transition at all Reynolds numbers studied, which carries about $96.2\%$ of the summed absolute contrasts at $Re=400$. The full exponent sequence is available at $N\le384$ at all three Reynolds numbers: at $\alpha=0$ the $n=4\to6$ contrast is $+0.00296\%$, $+0.0129\%$, and $+0.0103\%$, and the $n=6\to8$ contrast is $+0.00042\%$, $+6\times10^{-5}\%$, and $-0.0037\%$ at $Re=100$, $400$, and $1000$. Resolution is judged transition by transition against each contrast's own span. At $Re=1000$ both trailing steps are resolved, the $(6,8)$ step by a factor of $34$ with the same negative sign on all computed grids; at $Re=400$ that step lies inside its own span, and at $Re=100$ the $(4,6)$ step mixes Mach regimes on the coarse grids and is quoted as an upper bound. At $\alpha=0.1$ both trailing transitions are negative and individually resolved at $Re=400$ and $1000$. Increasing undulation amplitude weakens the primary contrast. The primary $\alpha=0$ contrast varies by at most $3.2\%$ of its own value over $N=128$--$768$ at each Reynolds number, while the absolute integrals retain roughly $1\%$ grid sensitivity at the production level and about $0.9\%$ between $N=384$ and $N=768$. Independent Taylor--Hood calculations reproduce the primary contrast to $0.002$ percentage points at $Re=400$ and $Re=1000$ on the finest meshes of each solver; the finite-element spread across its own meshes exceeds these gaps, so the comparison establishes consistency rather than convergence. Scope is limited to the tested superellipse family and a single undulation count; the unmatched Mach numbers of the coarse $n\ge6$ runs at $Re=100$ and the absence of $n=6,8$ fields at $N=768$ preclude universal scaling or grid-independent claims for the trailing transitions. Numerical uncertainties on $I_D(2,4)$ are $7.2\times10^{-4}$, $2.4\times10^{-3}$ and $1.4\times10^{-2}$ percentage points at $Re=100$, $400$ and $1000$, each measured over $N=128$--$768$; the spans over $N=128$--$384$ are $3.4\times10^{-4}$, $2.4\times10^{-3}$ and $1.3\times10^{-2}$~pp. The trailing contrasts have their own, generally smaller spans, collected in Table~SM.X.

Finally, the crossover of the local exponent ($\gamma_{\mathrm{loc}}=1$) near $Re\approx346$, observed in the exploratory $N=128$ Reynolds sweep where $Ma$ reaches $0.18$, is a descriptive feature of that coarse-grid trajectory and not a converged result; no conclusion depends on its precise location.

A Richardson extrapolation of $D^*$ and $P^*$ (Section~VII of the Supplementary Material) gives effective orders between $0.70$ and $1.07$ that are not stable under removal of a single grid, so no converged limit or grid-convergence index is claimed. The extrapolated contrast is also not robust: at $Re=100$ the four-level fit gives $+0.0233\%$, against $+0.0191\%$ with the three production grids and $+0.0257\%$ when the coarsest grid is dropped, a spread of a factor $1.35$. No extrapolated contrast is therefore quoted as a result. The dominant residual uncertainty is the lid-layer velocity gradient at coarse resolution: a second-order reconstruction changes $D^*$ by $-17\%$ at $N=128$, but by only $-0.24\%$ on the converged $N=768$ field, and an independent non-equilibrium MRT estimate agrees with the production finite-difference value to $0.16\%$ globally and $0.2\%$ on the lid row. The production (first-order) value is retained as the reported number.

Several targeted extensions would reduce the conservative spans reported here. A converged Richardson fit, with an effective order stable under removal of the coarsest grid, would replace the span by an extrapolated value with an explicit grid-convergence index; the present sequences do not provide such an order. Extending the $N=768$ series to $n=6$ and $8$ would give the trailing transitions the same resolution treatment as the primary one. A $Re=100$ Taylor--Hood calculation would test cross-solver agreement at the low-Reynolds end. Finally, a finer-grid Reynolds sweep would establish whether the crossover $Re\approx346$, currently a descriptive feature of the $N=128$ sweep, survives refinement.
